\begin{ouverture}[courbe]{5}{Équations et mise en équation}{Partie 1 · Nombres et calculs}
{Inconnue · Équations linéaires · Produits nuls · Validation}
\objectif{Résoudre une équation du premier degré.}
\objectif{Utiliser un produit nul ou un carré.}
\objectif{Mettre un problème en équation et valider sa solution.}
\prerequis{Calcul littéral et opérations inverses.}
\end{ouverture}
\pagestyle{otMain}\essentiel\fontsize{11}{15}\selectfont
\essSection{1}{Transformer une équation}
\begin{definition}Résoudre une équation, c’est trouver toutes les valeurs de l’inconnue qui rendent l’égalité vraie. Une solution se vérifie dans l’égalité initiale.\end{definition}
\begin{propriete}Ajouter ou soustraire une même quantité aux deux membres conserve les solutions. Multiplier ou diviser les deux membres par un même nombre non nul les conserve aussi.\end{propriete}
\begin{exemple}$5x-7=2x+8$ donne $3x-7=8$, puis $3x=15$, donc $x=5$. Contrôle : les deux membres initiaux valent 18.\end{exemple}
\begin{savoirfaire}[Organiser les étapes]Développer les parenthèses si nécessaire ; regrouper les termes en $x$ d’un côté et les constantes de l’autre ; isoler $x$ ; vérifier. Pour des fractions, multiplier tous les termes par un même dénominateur non nul.\end{savoirfaire}
\begin{attention}Une équation peut avoir une seule solution, aucune solution, ou tous les nombres pour solutions. $2x+3=2x+5$ est impossible ; $2(x+3)=2x+6$ est toujours vraie.\end{attention}
\clearpage
\essSection{2}{Produits, carrés et problèmes}
\begin{propriete}Un produit est nul si et seulement si au moins un de ses facteurs est nul. Ainsi $(x-4)(2x+3)=0$ donne $x=4$ ou $x=-3/2$.\end{propriete}
\begin{propriete}L’équation $x^2=a$ n’a pas de solution réelle si $a<0$. Si $a=0$, sa seule solution est 0. Si $a>0$, ses deux solutions sont $-\sqrt a$ et $\sqrt a$.\end{propriete}
\begin{savoirfaire}[Mettre en équation]Choisir une inconnue et préciser son sens et son unité. Exprimer les autres quantités en fonction de cette inconnue. Traduire l’égalité du problème. Résoudre puis revenir à la situation : un nombre d’objets est entier et une longueur est positive.\end{savoirfaire}
\begin{exemple}Pour deux tarifs $3x+12$ et $5x$, l’égalité se produit à $x=6$. Comparer ensuite des valeurs ou leur différence permet de répondre à une question de choix de tarif.\end{exemple}
\begin{attention}Diviser par $x$ peut perdre la solution $x=0$. Pour $x^2=5x$, écrire $x(x-5)=0$ et conserver les deux possibilités.\end{attention}
\clearpage\rubrique{Découverte}
\begin{blocLivre}{ex-5-1}\begin{exercice}[1]{Une opération inverse}
Résoudre $x+13=4$ et $-4x=18$.
\end{exercice}\end{blocLivre}
\begin{blocLivre}{ex-5-2}\begin{exercice}[1]{Deux étapes}
Résoudre $7x-5=30$.
\end{exercice}\end{blocLivre}
\begin{blocLivre}{ex-5-3}\begin{exercice}[1]{L’inconnue des deux côtés}
Résoudre $8x+3=5x-12$.
\end{exercice}\end{blocLivre}
\begin{blocLivre}{ex-5-4}\begin{exercice}[1]{Avec une parenthèse}
Résoudre $3(2x-5)=4x+9$.
\end{exercice}\end{blocLivre}
\begin{blocLivre}{ex-5-5}\begin{exercice}[1]{Une solution fractionnaire}
Résoudre $5x+2=9$ et donner la solution exacte.
\end{exercice}\end{blocLivre}

\clearpage\rubrique{Entraînement}
\begin{blocLivre}{ex-5-6}\begin{exercice}[2]{Une équation fractionnaire}
Résoudre $\dfrac{x}{3}+\dfrac{x}{4}=14$.
\end{exercice}\end{blocLivre}
\begin{blocLivre}{ex-5-7}\begin{exercice}[2]{Tout ou rien}
Résoudre $4(x-2)=4x-8$, puis $4(x-2)=4x+1$.
\end{exercice}\end{blocLivre}
\begin{blocLivre}{ex-5-8}\begin{exercice}[2]{Un produit nul}
Résoudre $(2x-7)(x+4)=0$.
\end{exercice}\end{blocLivre}
\begin{blocLivre}{ex-5-9}\begin{exercice}[2]{Une équation de carré}
Résoudre $x^2=49$, puis $x^2=12$, puis $x^2=-2$.
\end{exercice}\end{blocLivre}
\begin{blocLivre}{ex-5-10}\begin{exercice}[2]{Ne pas perdre zéro}
Résoudre $x^2-6x=0$. Expliquer pourquoi diviser directement par $x$ est risqué.
\end{exercice}\end{blocLivre}

\clearpage\rubrique{Problèmes}
\begin{blocLivre}{ex-5-11}\begin{exercice}[3]{Des billets de spectacle}
Un groupe achète 28 billets : les billets jeunes coûtent 6 €, les autres 10 €. La dépense totale est 232 €. Combien de billets jeunes ?
\end{exercice}\end{blocLivre}
\begin{blocLivre}{ex-5-12}\begin{exercice}[3]{Un rectangle à retrouver}
Un rectangle a un périmètre de 58 cm. Sa longueur dépasse sa largeur de 7 cm. Déterminer ses dimensions et son aire.
\label{fig-ex-5-12}\figureGeo{rectangle-equation}
\end{exercice}\end{blocLivre}
\begin{blocLivre}{ex-5-13}\begin{exercice}[3]{Retrouver le nombre initial}
On ajoute 4 à un nombre, on multiplie le résultat par 3, puis on soustrait le double du nombre de départ. Le résultat est 25. Retrouver le nombre.
\end{exercice}\end{blocLivre}
\begin{blocLivre}{ex-5-14}\begin{exercice}[3]{Un achat rentable}
Une bouteille filtrante coûte 38 €. Ses recharges coûtent 12 € par an. Elle évite un achat estimé à 30 € de bouteilles chaque année. À partir de quelle année entière la dépense cumulée devient-elle inférieure ou égale à l’ancien achat ?
\end{exercice}\end{blocLivre}
\begin{blocLivre}{ex-5-15}\begin{exercice}[3]{Deux aires égales}
Un carré de côté $x+2$ cm et un rectangle de côtés $x$ cm et $x+5$ cm ont la même aire. Les dimensions sont positives. Déterminer $x$ et l’aire commune.
\label{fig-ex-5-15}\figureGeo{aires-egales}
\end{exercice}\end{blocLivre}

\clearpage
